% ============================================================
% Chapter 8: Sustainability and Ethical Implications
%   Section: Equity and Bias Analysis
%     Subsection: Biases in the Processes
% Included from chapters/08-sustainability-ethics/03-equity-bias-analysis/equity-bias-analysis.tex
% ============================================================

\subsection{Biases in the Processes}

\begin{itemize}
    \item \textbf{The mission screen} encodes a view of what the
    publication does not cover --- politics, protests, crime, markets ---
    in nine written descriptions and their margins. It leaves out labour
    disputes and protests along with party politics, which is an
    editorial position, not a neutral filter, and it can leave out a
    constructive story whose words sit near one of those subjects (a
    prison-to-firefighter jobs programme was read as crime). The three
    positive-only outlets are exempt, which favours their editors'
    judgement over the screen's.
    \item \textbf{The ranking and the AI selection} encode a definition
    of positive impact written by the author. The ranking's weights were
    fitted on labels an AI assigned; the AI selection scores threat
    stories on topic highly and candidates without a summary low.
    \item \textbf{The search engines} decide what evidence exists for
    the system: their ranking, their refusals on a given day, and their
    coverage of Spanish sources. Google and Brave refused the
    metasearch engine during the pilot.
    \item \textbf{The models.} The embedding model places stories in a
    space learned from its training data; the sentiment model reads a
    problem being solved as negative, which is why it was kept out of
    the mission screen; and the verifier leans to \texttt{FALSE} on
    political claims whatever the evidence.
\end{itemize}
